\subsection{EQUIVALENCE CLASSES ; QUOTIENT SET}

Let \(f\) be a function, \(E\) its domain, \(F\) its graph. The relation `` \(x \in E\) and \(y \in E\) and \(f(x) = f(y)\) '' is an equivalence relation on \(E\) , called the equivalence relation associated with \(f\) . It is equivalent to the relation \((\exists z)((x, z) \in F \text{ and } (y, z) \in F)\) , i.e., to \((\exists z)((x, z) \in F \text{ and } (z, y) \in \overline{F})\) , and therefore its graph is \(\overline{F} \circ F\) .

We shall now show that every equivalence relation R on a set E is of this type. Let G be the graph of R. For each \(x \in E\) the (non-empty) set \(\mathrm{G}(x) \subset \mathrm{E}\) is called the equivalence class of x with respect to R; it is thus the set of all \(y \in E\) such that \(R\{x, y\}\) . Every set which can be written as \(\mathrm{G}(x)\) for some \(x \in E\) is called an equivalence class (with respect to R). An element of an equivalence class is called a representative of this class. The set of equivalence classes with respect to R (that is, the set of all objects of the form \(\mathrm{G}(x)\) for some \(x \in E\) ) is called the quotient set of E by R and is denoted by E/R. The mapping \(x \to \mathrm{G}(x)\) ( \(x \in E\) ) whose domain is E and whose target is E/R is called the canonical mapping of E onto E/R.

C55. Let R be an equivalence relation on a set E, and let p be the canonical mapping of E onto E/R. Then

\[
\mathrm{R} \left\{x, y \right\} \Longleftrightarrow (p (x) = p (y)).
\]

With the notation above, let x and y be elements of E such that

\[
(x, y) \in G.
\]

Then \(x \in E\) and \(y \in E\) ; let us show that \(G(x) = G(y)\) . Since \(y \in G(x)\) , we have (Proposition 1) \(G(y) \subset (G \circ G)(x) = G(x)\) . On the other hand, we also have \((y, x) \in G\) , whence \(G(x) \subset G(y)\) and therefore

\[
\mathbf {G} (x) = \mathbf {G} (y),
\]

i.e., \(p(x) = p(y)\) . Conversely, if \(G(x) = G(y)\) , we have \(y \in G(y) = G(x)\) , so that \((x, y) \in G\) . This completes the proof.

A section of the canonical mapping \(p\) of \(\mathbf{E}\) onto \(\mathbf{E} / \mathbf{R}\) (§ 3, no. 8, Definition 11) is called more briefly a section of \(\mathbf{E}\) (with respect to the relation \(\mathbf{R}\) ).

\minorheading{Examples}

\begin{enumerate}
\def\labelenumi{(\arabic{enumi})}
\item
  Let R be the equivalence relation `` \(x \in E\) and \(y \in E\) and \(x = y\) '' on a set E. The equivalence class of \(x \in E\) is then the set \(\{x\}\) , and the canonical mapping \(x \to \{x\}\) of E onto E/R is bijective.
\item
  Let E, F be two sets and let R be the equivalence relation on \(E \times F\) associated with the mapping \(pr_{1}\) of \(E \times F\) onto E. The equivalence classes with respect to R are the sets of the form \(\{x\} \times F\) , where \(x \in E\) ; the mapping \(x \to \{x\} \times F\) is a bijection of E onto \((\mathbf{E} \times \mathbf{F}) / R\) .
\end{enumerate}

Let R be an equivalence relation on a set E. The quotient set E/R is a subset of \(\mathfrak{P}(E)\) , and the identity mapping of E/R is a partition of E ( \(\S4\) , no. 7); for if G is the graph of R, we have \(x \in \mathrm{G}(x)\) for all \(x \in E\) , and if two equivalence classes \(\mathbf{G}(x)\) and \(\mathbf{G}(y)\) have a common element \(z\) , then \(\mathbf{R}\{x, z\}\) and \(\mathbf{R}\{y, z\}\) , so that \(\mathbf{G}(x) = \mathbf{G}(y)\) . Furthermore, the relation

\[
(\exists \mathrm{X}) (\mathrm{X} \in \mathrm{E} / \mathrm{R} \text { and } x \in \mathrm{X} \text { and } y \in \mathrm{X})
\]

is equivalent to R\{x, y\}.

\(\mathbb{J}\) Conversely, let \((X_{i})_{i\in I}\) be a partition of a set E. It is immediately verified that the relation \((\exists i)(i\in I \text{ and } x\in X_{i} \text{ and } y\in Y_{i})\) is an equivalence relation R on E; the equivalence classes with respect to R are just the sets \(X_{i}\) of the partition, and the mapping \(i\to X_{i}\) is a bijection of I onto E/R. Every subset S of E such that, for each \(i\in I\) , the set \(S\cap X_{i}\) consists of a single element is called a system of representatives (or a transversal) of the equivalence classes with respect to R. This name is also used to denote any injection of a set K into E such that the image of K under this injection is a system of representatives of the equivalence classes with respect to R; for example, any section of E with respect to R.

